\documentclass{article}

% if you need to pass options to natbib, use, e.g.:
%     \PassOptionsToPackage{numbers, compress}{natbib}
% before loading neurips_2020

% ready for submission
% \usepackage{neurips_2020}

% to compile a preprint version, e.g., for submission to arXiv, add add the
% [preprint] option:
%     \usepackage[preprint]{neurips_2020}

% to compile a camera-ready version, add the [final] option, e.g.:
    \usepackage[final]{project}

% to avoid loading the natbib package, add option nonatbib:
     % \usepackage[nonatbib]{neurips_2020}

\usepackage[utf8]{inputenc} % allow utf-8 input
\usepackage[T1]{fontenc}    % use 8-bit T1 fonts
\usepackage{hyperref}       % hyperlinks
\usepackage{url}            % simple URL typesetting
\usepackage{booktabs}       % professional-quality tables
\usepackage{amsfonts}       % blackboard math symbols
\usepackage{nicefrac}       % compact symbols for 1/2, etc.
\usepackage{microtype}      % microtypography
\usepackage{caption}
\usepackage{subcaption}
\usepackage{amsmath}
\usepackage{amsfonts}
\usepackage{amssymb}
\usepackage{graphicx}

\title{Learning When to Think: Reinforcement Learning for Adaptive Reasoning Computation in LLMs}

\author{%
  Ivan Andhika \\
  School of Computing\\
  University of Utah\\
  \texttt{u1590903@utah.edu} \\
  \And
  Hao Ren \\
  School of Computing\\
  University of Utah\\
  \texttt{u1527543@utah.edu} \\
  \And
  Ammon Gleason \\
  School of Computing\\
  University of Utah\\
  \texttt{u1221580@utah.edu} \\
}

\begin{document}

\maketitle

\begin{abstract}
Large language models (LLMs) can improve reasoning performance by increasing test-time computation---generating longer chains of reasoning or sampling multiple candidate solutions. However, most approaches allocate a \emph{fixed} amount of computation for every input, which is inefficient: simple problems receive unnecessary reasoning while difficult problems may require additional exploration. We investigate whether reinforcement learning (RL) can train LLMs to \emph{adaptively} allocate reasoning computation during inference. We formulate reasoning as a sequential decision-making problem in which the model chooses among actions---continuing reasoning, verifying an intermediate result, sampling an alternative solution, or terminating---and learn a policy that dynamically determines how much effort each problem receives. Training uses RL with a reward that balances answer correctness against computational cost. We evaluate on GSM8K using Qwen2.5-0.5B with LoRA fine-tuning, comparing the learned adaptive policy against fixed-compute baselines including standard chain-of-thought and self-consistency sampling. We assess both task accuracy and compute efficiency, reporting metrics such as reasoning token usage and cost-per-correct-answer across problem difficulty levels.
\end{abstract}

\section{Introduction}

Recent work has demonstrated that increasing the amount of computation a language model expends at inference time can substantially improve its reasoning capabilities~\cite{wei2022chainofthought, snell2025scaling}. Techniques such as chain-of-thought (CoT) prompting~\cite{wei2022chainofthought}, self-consistency decoding~\cite{wang2023selfconsistency}, and iterative self-refinement~\cite{kumar2025training} have shown that \emph{thinking more} often leads to better answers on tasks requiring multi-step reasoning.

However, these methods typically apply a uniform computation budget to every input. A one-step arithmetic problem receives the same lengthy chain-of-thought treatment as a complex eight-step word problem, and self-consistency always samples a fixed number $k$ of candidate solutions regardless of difficulty. This rigidity wastes computation on easy inputs and may under-allocate effort on hard ones. The problem is especially acute for smaller models deployed under resource constraints, where every token of wasted computation has a meaningful cost.

This inefficiency motivates a natural question: \emph{can a language model learn to decide, for each problem, how much reasoning effort to invest?} We frame this as a reinforcement learning problem. At each step of its reasoning process, the model acts as an agent that selects from a discrete set of actions: \textbf{continue} reasoning, \textbf{verify} the current intermediate result, try an \textbf{alternative} solution path, or \textbf{terminate} and output a final answer. The reward signal balances correctness against token cost, incentivizing the model to be both accurate and efficient.

This problem is directly relevant to the themes of this course (CS 5955/6955: Advanced AI). It connects reinforcement learning---particularly policy optimization in sequential decision-making---with modern large language models. The adaptive computation question also relates to exploration--exploitation trade-offs studied in multi-armed bandits: the model must decide whether to ``exploit'' its current reasoning (terminate) or ``explore'' further (continue, verify, or try an alternative).

Our contributions are: (1) a formulation of adaptive test-time computation as an RL problem with a structured action space; (2) an empirical comparison of training approaches including supervised fine-tuning (SFT), Direct Preference Optimization (DPO), and Group Relative Policy Optimization (GRPO); and (3) a difficulty-aware evaluation framework that analyzes whether the learned policy allocates more computation to harder problems.

\section{Related Work}

\textbf{Chain-of-thought reasoning.} Wei et al.~\cite{wei2022chainofthought} demonstrated that prompting LLMs to produce intermediate reasoning steps dramatically improves performance on arithmetic, commonsense, and symbolic reasoning tasks. Wang et al.~\cite{wang2023selfconsistency} extended this with self-consistency, which samples multiple reasoning paths and selects the answer by majority vote. Both methods use a fixed computation budget per problem. Our work aims to make this budget adaptive.

\textbf{Test-time compute scaling.} Snell et al.~\cite{snell2025scaling} showed that optimally allocating test-time compute on a per-problem basis can be more effective than scaling model parameters, with a smaller model using adaptive compute outperforming a 14$\times$ larger model. Their work provides theoretical motivation for our approach but focuses on search-based strategies rather than learning an explicit policy via RL.

\textbf{RL for LLM reasoning.} DeepSeek-R1~\cite{deepseek2025r1} demonstrated that reinforcement learning alone---without supervised reasoning traces---can incentivize emergent reasoning behaviors including self-reflection and verification. DeepSeekMath~\cite{shao2024deepseekmath} introduced Group Relative Policy Optimization (GRPO), which computes advantages relative to a group of rollouts rather than requiring a separate critic model, reducing memory overhead. We adopt GRPO as one of our training methods. Kumar et al.~\cite{kumar2025training} proposed SCoRe, a multi-turn RL method for training self-correction, showing that RL can teach models to iteratively improve their own outputs.

\textbf{Adaptive reasoning depth.} Most closely related to our work, Fang et al.~\cite{fang2025thinkless} proposed \emph{Thinkless}, which trains an LLM to choose between short-form and extended thinking using two control tokens and a Decoupled GRPO algorithm, reducing extended thinking by 50--90\%. Xiang et al.~\cite{xiang2025justenough} introduced Adaptive Length Penalties (ALP) that scale inversely with per-prompt solve rate, cutting token usage by $\sim$50\%. Our approach differs from Thinkless in that we define a \emph{richer action space} with four distinct actions (continue, verify, alternative, terminate) rather than a binary short/long decision, enabling finer-grained control over the reasoning process.

\textbf{Process supervision.} Lightman et al.~\cite{lightman2024lets} showed that providing step-level feedback (process supervision) significantly outperforms outcome-only supervision for training reward models on mathematical reasoning. This motivates our inclusion of a \emph{verify} action that allows the model to pause and evaluate intermediate steps.

\textbf{Reasoning and acting.} Yao et al.~\cite{yao2023react} proposed ReAct, which interleaves reasoning traces with task-specific actions. Our formulation shares the insight that reasoning should be an active, decision-driven process, but we focus specifically on controlling the \emph{depth} and \emph{style} of reasoning rather than interfacing with external tools.

\section{Method}

\subsection{Problem Formulation}

We model the reasoning process as a Markov Decision Process (MDP). Given a problem $q$, the agent (LLM) generates a sequence of reasoning steps. At each step $t$, the state $s_t$ consists of the problem text and all reasoning generated so far. The agent selects an action $a_t$ from the action space $\mathcal{A} = \{\texttt{continue}, \texttt{verify}, \texttt{alternative}, \texttt{terminate}\}$:

\begin{itemize}
    \item \textbf{Continue}: Extend the current chain of reasoning by generating an additional reasoning step.
    \item \textbf{Verify}: Prompt the model to check its most recent intermediate result before proceeding.
    \item \textbf{Alternative}: Discard the current reasoning path and sample a fresh solution attempt.
    \item \textbf{Terminate}: Stop reasoning and extract the final answer from the current context.
\end{itemize}

The episode ends when the agent selects \texttt{terminate} or reaches a maximum step limit $T$.

\subsection{Reward Function}

Upon termination, the agent receives a reward that balances correctness and efficiency:
\begin{equation}
    R = \underbrace{r_{\text{acc}}}_{\text{correctness}} + \underbrace{r_{\text{fmt}}}_{\text{format bonus}} - \underbrace{\lambda \cdot \frac{n_{\text{tokens}}}{n_{\text{max}}}}_{\text{cost penalty (if correct)}}
\end{equation}
where $r_{\text{acc}} = 1$ if the predicted answer matches the gold answer (within tolerance) and $0$ otherwise, $r_{\text{fmt}} = 0.1$ if the output follows the expected answer format, $\lambda$ is a hyperparameter controlling the accuracy--efficiency trade-off, and $n_{\text{tokens}}$ is the total number of tokens generated. Crucially, the cost penalty is applied \emph{only when the answer is correct}, so the model is not penalized for exploring on problems it cannot yet solve.

\subsection{Training Approaches}

We investigate three training methods, each building on the previous:

\textbf{Supervised Fine-Tuning (SFT).} We first distill reasoning capabilities from a stronger teacher model (Gemini 2.5 Flash) or from GSM8K's gold step-by-step solutions. The student model (Qwen2.5-0.5B-Instruct) is fine-tuned on these traces using LoRA (rank 16) on the attention projection layers.

\textbf{Direct Preference Optimization (DPO).} After SFT, we generate multiple candidate solutions per problem. Correct answers form ``chosen'' examples and incorrect answers form ``rejected'' examples. DPO~\cite{rafailov2023direct} trains the model to prefer correct solutions without requiring an explicit reward model.

\textbf{Group Relative Policy Optimization (GRPO).} Following Shao et al.~\cite{shao2024deepseekmath}, we generate $K$ rollouts per problem, score each with the reward function, compute group-normalized advantages, and update the policy to increase the probability of high-reward rollouts. GRPO avoids the memory cost of a separate critic model, making it feasible for our compute budget.

\section{Experimental Design}

\subsection{Hypotheses}

We test the following hypotheses:

\begin{enumerate}
    \item \textbf{H1 (Adaptive compute)}: An RL-trained adaptive policy will achieve a better accuracy--efficiency trade-off than fixed-compute baselines (CoT, self-consistency) on GSM8K, measured by cost-per-correct-answer.
    \item \textbf{H2 (Difficulty sensitivity)}: The adaptive policy will learn to allocate more tokens to harder problems and fewer tokens to easier problems, as measured by a positive correlation between ground-truth problem difficulty and average tokens generated.
    \item \textbf{H3 (Action specialization)}: The distribution of selected actions will vary with problem difficulty---the model will \texttt{terminate} earlier on easy problems and use \texttt{verify}/\texttt{alternative} more frequently on hard problems.
\end{enumerate}

\subsection{Domain and Data}

All experiments use the GSM8K benchmark~\cite{cobbe2021gsm8k}, which contains 8,792 grade-school math word problems with human-written step-by-step solutions. Problems range from 1-step arithmetic to 8+ step multi-part reasoning, providing natural variation in difficulty. We classify problems by counting calculation steps (annotated as \texttt{<<expr=result>>}) in the gold solutions and binning into \emph{easy} ($\leq 2$ steps), \emph{medium} (3--4 steps), and \emph{hard} ($\geq 5$ steps).

\subsection{Experimental Setup}

\textbf{Base model:} Qwen2.5-0.5B-Instruct with LoRA fine-tuning (rank 16, $\alpha=32$) on attention layers.

\textbf{Baselines:}
\begin{itemize}
    \item \textbf{CoT baseline}: Single-pass chain-of-thought generation (fixed 512 token budget).
    \item \textbf{Self-consistency ($k=5$)}: Sample 5 independent CoT solutions and take the majority vote.
\end{itemize}

\textbf{Our methods:}
\begin{itemize}
    \item \textbf{SFT + DPO}: Distillation followed by preference optimization.
    \item \textbf{GRPO}: Direct RL training with group-relative advantages and KL regularization.
\end{itemize}

\subsection{Evaluation Metrics}

\begin{itemize}
    \item \textbf{Accuracy}: Fraction of problems answered correctly.
    \item \textbf{Average tokens}: Mean tokens generated per problem.
    \item \textbf{Cost-per-correct}: Total tokens divided by number of correct answers (lower is more efficient).
    \item \textbf{Difficulty-conditioned metrics}: All of the above broken down by easy/medium/hard bins.
    \item \textbf{Action distribution}: Frequency of each action type by difficulty level.
\end{itemize}

We evaluate on 200 problems from the GSM8K test set and report 95\% Wilson score confidence intervals for accuracy. Ablations over $\lambda \in \{0.01, 0.05, 0.1, 0.2\}$ and maximum steps $T \in \{3, 5, 7\}$ will characterize the sensitivity of the accuracy--efficiency trade-off.

\bibliographystyle{plain}
\bibliography{sample}

\end{document}
